\documentclass[10pt,a4paper]{amsart}
\usepackage[margin=19mm]{geometry}
\usepackage{amsmath,amssymb,mathtools}
\usepackage[scr=rsfs,cal=euler]{mathalfa}
\usepackage{xfrac}
\usepackage[unicode,hidelinks,bookmarks=false]{hyperref}
\newcommand{\ie}{i.e.\ }
\newcommand{\cf}{cf.\ }
\newcommand{\resp}{respectively\ }
\newcommand{\Subsection}{\subsection}
\providecommand{\nobreakdash}{\nobreak\hbox{-}}
\setlength{\emergencystretch}{2em}
\multlinegap0pt
\usepackage{bm}
\usepackage{bbm}
\usepackage{dsfont}
\usepackage{xspace}
\usepackage{cancel}
\usepackage{mathtools}
\usepackage{amsthm}
\makeatletter
\@ifundefined{define}{\newtheorem{define}{Define}}{}
\@ifundefined{claim}{\newtheorem{claim}[define]{Claim}}{}
\@ifundefined{lemma}{\newtheorem{lemma}[define]{Lemma}}{}
\makeatother
\newcommand{\N}{\mathbb{N}}
\newcommand{\Z}{\mathbb{Z}}
\title[Non-rectifiable Delone sets under pointwise co-Lipschitz bij]{Non-rectifiable Delone sets under pointwise co-Lipschitz bijections}
\author{Ashwin Bhat, Michael Dymond}
\date{Source: arXiv:2608.17637v1 (CC BY 4.0)}
\begin{document}
\maketitle

\noindent\emph{Source and licence.} Non-rectifiable Delone sets under pointwise co-Lipschitz bijections. Version arXiv:2608.17637v1 (CC BY 4.0), distributed under the Creative Commons Attribution 4.0 International licence, as recorded in the arXiv licence metadata for this exact version and shown on the abstract page https://arxiv.org/abs/2608.17637v1. Full rights notice: licenses/arxiv-2608.17637-CC-BY-4.0.txt.\par\smallskip

\noindent\textit{Editorial scope.} Counted unit: the Lemma whose source label is \texttt{lemma: cohom examples} in the article (source0.tex lines 1044--1046, source label \texttt{lemma: cohom examples}), delivered together with the article's own complete original proof of it (source0.tex lines 1047--1153, one \texttt{proof} environment of 107 lines). No mathematics is written, repaired or re-derived: statement and proof are the article's own text line for line. Required context retained uncounted: eq: definition of f (lines 905--907); lemma: 5-lipschitz (lines 935--937); lemma: bilipschitz properties (lines 973--985); enumerate: gd (lines 979--984). Every definition, display and label of those lines is retained unchanged. Omitted from this package and not redistributed: 1--904, 908--934, 938--972, 985--1043, 1154--1273. Nothing inside a retained excerpt is omitted or altered: every retained body is delivered exactly as the article's lines are written, so no omission receipt of any kind accompanies this package. Citations: the retained excerpts carry no citation command at all, so no bibliography entry of the article is transcribed and no reference is redistributed. Reference adaptation: none was needed and none was made; every reference inside the delivered unit resolves inside it, so no unresolved reference or citation is produced. Printed slips kept literal: no printed word, symbol, hypothesis or number of the counted statement or of its proof was altered. Layout and class: the article's own class is \texttt{article} with packages including graphicx, amsmath, amssymb, enumerate, amsthm, todonotes, mathtools, babel, fontenc, inputenc, lmodern, tikz, thmtools, thm-restate; the wrapper is the lane's pinned \texttt{amsart} editorial wrapper at 10pt on A4, which restates the theorem-like environments the retained text needs and adds the article's own macro definitions that the retained text needs (\texttt{\textbackslash N}, \texttt{\textbackslash Z}), with extra wrapper packages (bm, bbm, dsfont, xspace, cancel, mathtools). The delivered numbering is the unit's own. Assets: no retained span contains a figure environment, a graphics-inclusion command, a PDF-page inclusion, a table or a TikZ picture; no image file is copied into this package, the package declares no asset dependency and no image file is redistributed with it. Licence: version arXiv:2608.17637v1 of this article is distributed under the Creative Commons Attribution 4.0 International licence, as recorded in the arXiv licence metadata for this exact version and shown on the abstract page https://arxiv.org/abs/2608.17637v1. A full rights notice is delivered at licenses/arxiv-2608.17637-CC-BY-4.0.txt. Characters: every retained excerpt is pure ASCII, so the delivered engine, XeLaTeX with its default Latin Modern text font, reports no missing character.
\begin{equation}\label{eq: definition of f}
    g(0) = 0, \qquad g|_{\N} = g_+ \qquad \text{and} \qquad g(x) = -g_+(-x) \qquad \text{for all $x \in -\N$}.
\end{equation}

\begin{lemma} \label{lemma: 5-lipschitz}
    For any sequence $(N_i)_{i \in \N} \subseteq \N$, the bijection $g:\Z \to \Z$ given by \eqref{eq: definition of f} is $5-$Lipschitz.
\end{lemma}

\begin{lemma} \label{lemma: bilipschitz properties}
    Let $d \in \N$, $(N_i)_{i \in \N} \subseteq \N$, the bijection $g:\Z \to \Z$ given by \eqref{eq: definition of f}, dependent on $(N_i)_{i \in \N}$, and the bijection $g^{(d)}:\Z^d \to \Z^d$ given by
    \begin{equation*} \label{eq: fd}
     g^{(d)}(x_1,...,x_d) = (g(x_1),x_2,...,x_d).
 \end{equation*}
    Then
    \begin{enumerate}
    \item $g^{(d)}$ is $5-$Lipschitz;
    \item \label{enumerate: gd} $\max\{||g^{(d)}(x)-x||,||(g^{(d)})^{-1}(x)-x||\} \leq 3\max_{j \in [i]}N_j \text{ whenever $||x|| \leq 3\sum_{j=1}^iN_j$}$ for some $i \in \N;$
        \item \label{point 3, lemma 4.9} $\left(N_{i+1}(\sum_{j=1}^{i}N_j)^{-1}\right)_{i \in \N}$ is bounded if and only if $g^{(d)}$ is pointwise co-Lipschitz;
        \item \label{point 4} $(N_i)_{i \in \N}$ is bounded if and only if $g^{(d)}$ is bi-Lipschitz.
    \end{enumerate}
\end{lemma}

\begin{lemma} \label{lemma: cohom examples}
Let $d \in \N$ and $\omega:[0,2) \to [0,\infty)$ be a concave, strictly increasing function with $\omega(0) = 0$. Then there exists a Lipschitz, $\omega-$co-homogeneous self-bijection of $\Z^d$ which is not $\widetilde{\omega}-$co-homogeneous for any concave, increasing function $\widetilde{\omega}$ satisfying $\widetilde{\omega}(0)=0$ and $\lim_{t \to 0^{+}}\tfrac{\widetilde{\omega}(t)}{\omega(t)}=0$. 
\end{lemma}

\begin{proof}
Without loss of generality, we may assume $\omega(1) = 1$, since a function is $\omega-$co-homogeneous if and only if it is $\frac{\omega}{\omega(1)}-$co-homogeneous. Let the sequence of natural numbers $(N_i)_{i \in \N}$ given by
$$N_i = \left\lceil \tfrac{1}{\Psi^{-1}(i)}-\tfrac{1}{\Psi^{-1}(i-1)}\right\rceil \qquad \text{with } \Psi(u)\coloneqq \int_{u}^{1}\frac{dt}{t \omega(t)}, \qquad u \in (0,1]$$
and let $g^{(d)}:\Z^d \to \Z^d$ be the Lipschitz bijection  dependent on $(N_i)_{i \in \N}$ from Lemma~\ref{lemma: bilipschitz properties}.
We will show $g^{(d)}$ is $\omega-$co-homogeneous. 
    First, we justify that $\Psi:(0,1] \to [0,\infty)$ is well-defined, invertible and surjective, and that $(N_i)_{i \in \N}$ is indeed a sequence of natural numbers. Since $\omega$ is an increasing function, we have that $\omega(t) \leq \omega(1) = 1$ for each $t \in (0,1].$ Moreover, $\omega$ is concave and thus continuous on $(0,1]$. Hence, by the comparison test for integrals $$\int_u^1 \frac{dt}{t\omega(t)} \leq \int_u^1 \frac{dt}{t} = -\log u, \qquad u \in (0,1],$$ so we have that $\Psi$ is finite and thus well-defined on $(0,1]$. Next, since $\omega(t) \leq 1$ for each $t \in (0,1],$ we have that
    $$\Psi(u)-\Psi(u+h) = \int_u^{u+h} \frac{dt}{t\omega(t)} \geq h\inf_{t \in [u,u+h]}\frac{1}{t\omega(t)} \geq h>0, \qquad u \in [0,1),\;h\in (0,1-u]$$
    so $\Psi$ is strictly decreasing, and hence $\Psi$ is invertible. We also have by the concavity of $\omega$ that $\omega(t) \geq t\omega(1)=t$ for each $t \in (0,1],$ so by the comparison test for integrals again
    $$\Psi(u) =\int_u^1 \frac{dt}{t\omega(t)} \geq \int_u^1 \frac{dt}{t^2} = \frac{1}{u}-1 \xrightarrow{u \to 0^+}\infty.$$ Hence, since $\Psi$ is differentiable and thus continuous by the fundamental theorem of calculus, then by the intermediate value theorem we have that $\Psi$ is surjective.
    Finally, since $\Psi$ is strictly decreasing, we have that $\Psi^{-1}$ is strictly decreasing too, so $N_i$ is a natural number for each $i \in \N.$

    Now we show $g^{(d)}$ is $\omega-$co-homogeneous. First, for each $i \in [0,\infty)$ we write $S_i \coloneqq \frac{1}{\Psi^{-1}(i)}.$ Fixing $i \in \N$
    \begin{equation*}
        1 = i-(i-1) = \Psi\left(\tfrac{1}{S_{i}}\right) - \Psi\left(\tfrac{1}{S_{i-1}}\right) = \int_{\frac{1}{S_i}}^{\frac{1}{S_{i-1}}}\frac{dt}{t\omega(t)}.
    \end{equation*}
    Now since $\frac{1}{t\omega(t)}$ is decreasing in $t$, and $S_i \geq S_{i-1}$ since $\Psi^{-1}$ is decreasing, we can bound the right-hand integral to obtain
    \begin{equation*}
        \left(\tfrac{1}{S_{i-1}}-\tfrac{1}{S_i}\right) \frac{1}{\frac{1}{S_i}\omega(\frac{1}{S_i})} \leq   1 \leq  \left(\tfrac{1}{S_{i-1}}-\tfrac{1}{S_i}\right) \frac{1}{\frac{1}{S_{i-1}}\omega(\frac{1}{S_{i-1}})},
    \end{equation*} so 
$$\omega\left(\tfrac{1}{S_i}\right) \geq S_i\left(\tfrac{1}{S_{i-1}}-\tfrac{1}{S_i}\right)  = \tfrac{S_i-S_{i-1}}{S_{i-1}}$$ and
$$\omega\left(\tfrac{1}{S_{i-1}}\right) \leq S_{i-1}\left(\tfrac{1}{S_{i-1}}-\tfrac{1}{S_i}\right) = \tfrac{S_i-S_{i-1}}{S_i}.$$
    Hence, we obtain 
\begin{equation}\label{eq: Si}
        N_i = \left\lceil S_i-S_{i-1}\right\rceil \leq \left\lceil S_{i-1}\omega\left(\tfrac{1}{S_i}\right)\right\rceil \leq 1+ S_{i-1}\omega\left(\tfrac{1}{S_i}\right)
    \end{equation} and
    \begin{equation}\label{eq: Si 2}
        N_i = \left\lceil S_i-S_{i-1}\right\rceil \geq  S_i-S_{i-1} \geq S_{i}\omega\left(\tfrac{1}{S_{i-1}}\right).
    \end{equation}
    \begin{claim}\label{claim: I}
       For each $i \in \N$ we have that
        $S_i \geq i.$
    \end{claim}
    \begin{proof}
Notice for each $i \in \N$ that
\begin{equation} \label{eq: eq in claim}
    \Psi\left(\tfrac{1}{i}\right) = \int_{\frac{1}{i}}^1 \frac{dt}{t\omega(t)} \leq \int_{\frac{1}{i}}^1 \frac{dt}{t^2} = i-1\leq i,
\end{equation}
 where for the first inequality we used the concavity of $\omega$ to say $\omega(t) \geq \omega(1)t = t$ for each $t \in (0,1].$ Since $\Psi$ is decreasing, applying $\Psi^{-1}$ to both sides of \eqref{eq: eq in claim} gives for each $i \in \N$ that $i\leq S_i.$
    \end{proof}
    By Claim \ref{claim: I}, for each $i\in \N$ we have that
  \begin{equation} \label{eq: sum Nj condition 2}
        \sum_{j=1}^iN_j = \sum_{j=1}^i\lceil S_j-S_{j-1}\rceil \leq \sum_{j=1}^i (S_j-S_{j-1}+1) = S_i -S_0+i\leq S_i+i \leq 2S_i.
 \end{equation} Similarly, since $S_0 = 1,$ notice
 \begin{equation} \label{eq: sum Nj condition prime}
        \sum_{j=1}^iN_j = \sum_{j=1}^i\lceil S_j-S_{j-1}\rceil \geq \sum_{j=1}^i S_j-S_{j-1} = S_i -S_0
 = S_i -1.
 \end{equation} Note that the last expression of \eqref{eq: sum Nj condition prime} is at least $\frac{S_i}{2}$ whenever $S_i \geq 2.$ By Claim \ref{claim: I}, $S_i \geq i \geq 2$ for each $i \geq 2.$ If $i = 1$ and $S_1 =1,$ the left-hand expression of \eqref{eq: sum Nj condition prime} gives $N_1$ which is at least $1$ by definition, and so $N_1 \geq \frac{S_1}{2}.$ Together, we obtain for each $i \in \N$ that
 \begin{equation} \label{eq: sum Nj condition}
     \sum_{j=1}^i N_j \geq \tfrac{S_i}{2}.
 \end{equation}
  Now fix $R>0.$ We have that $R \in (3\sum_{j=1}^{i-1}N_j,3\sum_{j=1}^{i}N_j]$ for some $i \in \N.$ Thus, using \eqref{eq: Si}, \eqref{eq: sum Nj condition 2} and \eqref{eq: sum Nj condition}
    \begin{equation} \label{eq: 4max}
        \begin{split}
            N_i &\leq 1+S_{i-1}\omega\left(\tfrac{1}{S_i}\right) \leq 1+2\left(\sum_{j=1}^{i-1}N_j\right)\omega\left(\frac{2}{\sum_{j=1}^iN_j}\right)
             \\&\leq 1+\tfrac{2R}{3}\omega\left(\tfrac{6}{R}\right) \leq R\omega\left(\tfrac{1}{R}\right) +\tfrac{2R}{3}6\omega\left(\tfrac{1}{R}\right) = 5R\omega\left(\tfrac{1}{R}\right),
        \end{split}
    \end{equation} where in the fourth inequality we used the concavity of $\omega$ to show ${\omega(\frac{6}{R}) \leq 6\omega(\frac{1}{R})}$ and $1=\omega(1) \leq R\omega(\frac{1}{R}).$
    Now we need the following claim.
    \begin{claim} \label{claim: nondec}
        $(N_i)_{i \in \N}$ is a non-decreasing sequence.
    \end{claim}
    \begin{proof}
    First, we claim that $\frac{\omega(t)}{t}$ is a non-increasing function on $(0,1]$. Let $t_1,t_2 \in (0,1]$ with $t_1\leq t_2$. Since $\omega$ is concave and $\omega(0) = 0$
        $$\omega(t_1) =\omega\left(\left(1-\tfrac{t_1}{t_2}\right)0+\tfrac{t_1}{t_2} t_2\right) \geq \left(1-\tfrac{t_1}{t_2}\right)\omega(0) + \tfrac{t_1}{t_2}\omega(t_2) = \tfrac{t_1}{t_2}\omega(t_2),$$
        so $\frac{\omega(t_1)}{t_1} \geq \frac{\omega(t_2)}{t_2},$ as required.
        Now by the fundamental theorem of calculus, we have that
        $$\Psi'(u) = -\tfrac{1}{u\omega (u)}, \qquad u \in (0,1].$$ Hence, for each $i \in [0,\infty)$ we have that
        \begin{equation*}
            (\Psi^{-1}(i))' = \tfrac{1}{\Psi'(\Psi^{-1}(i))} = -\Psi^{-1}(i)\omega(\Psi^{-1}(i)).
        \end{equation*}
        Thus, for each $i \in [0,\infty)$
        \begin{equation*}
        \begin{split}
            S_i' &= \left(\tfrac{1}{\Psi^{-1}(i)}\right)' = -\tfrac{1}{(\Psi^{-1}(i))^2} \left(\Psi^{-1}(i)\right)' = -\tfrac{1}{(\Psi^{-1}(i))^2} \cdot-\Psi^{-1}(i)\omega(\Psi^{-1}(i)) = \tfrac{\omega(\Psi^{-1}(i))}{\Psi^{-1}(i)}.
            \end{split}
        \end{equation*} Since both $\Psi^{-1}$ and $\frac{\omega(t)}{t}$ are non-increasing functions, we have that $S_i'$ is a non-decreasing function, and so $S_i$ is convex.
        Therefore, for each $i \in \N$
        $$2S_i = 2S_{\frac{1}{2}(i-1)+\frac{1}{2}(i+1)} \leq 2\left({\tfrac{1}{2}S_{i-1} +\tfrac{1}{2}S_{i+1}}\right)=S_{i-1}+S_{i+1}$$
        so we have that
        $$S_{i+1}-S_i \geq S_i - S_{i-1} \qquad \text{for each } i \in \N.$$ 
        Taking the ceiling of both sides, we obtain that $N_{i+1} \geq N_i$ for each $i \in \N,$ as required.
    \end{proof}
    With the insertion of Claim \ref{claim: nondec}, point \ref{enumerate: gd} of Lemma \ref{lemma: 5-lipschitz} gives that
    $$||(g^{(d)})^{-1}(x)-x|| \leq 3N_i \qquad \text{whenever $||x|| \leq 3\sum_{j=1}^iN_j$}.$$
    Then fixing $x,y \in B(0,R) \cap \Z^d$ with $x \neq y,$ we have that
    \begin{equation*}
    \begin{split}
        ||(g^{(d)})^{-1}(x) - (g^{(d)})^{-1}(y)|| &\leq ||(g^{(d)})^{-1}(x)-x|| + ||(g^{(d)})^{-1}(y)-y|| + ||x-y|| \leq 6N_i + ||x-y||\\
        &\leq 30R\omega\left(\tfrac{1}{R}\right)+R\omega\left(\tfrac{||x-y||}{R}\right)\leq 31R\omega\left(\tfrac{||x-y||}{R}\right),
         \end{split}
    \end{equation*} where for the penultimate inequality we used \eqref{eq: 4max} and that $\omega(t) \geq t$ for each $t \in (0,1]$ by the concavity of $\omega$, and for the last inequality we used that $||x-y|| \geq 1.$ Thus, $g^{(d)}$ is $\omega-$co-homogeneous.

    It remains to check $g^{(d)}$ does not witness any `better' co-homogeneity. Fix a concave, strictly increasing function $\widetilde{\omega}:[0,2) \to [0,\infty)$ with $\widetilde{\omega}(0) = 0$ and $\lim_{t \to 0^+} \frac{\widetilde{\omega}(t)}{\omega(t)} = 0$. 
    Define the sequences $(R_i)_{i \in \N}\subseteq \N,$ and $(x_i)_{i \in \N},(y_i)_{i \in \N} \subseteq \N^d$ as
    $$R_i = 3\sum_{j=1}^{i+1}N_j, \qquad x_i = \left(1+3\sum_{j=1}^{i}N_j,0,...,0\right),\qquad y_i =\left(-2+3\sum_{j=1}^{i}N_j,0,...,0\right) .$$
  Notice for each $i \in \N$ that $x_i,y_i \in B(0,R_i) \cap \Z^d.$ 
  By \eqref{eq: sum Nj condition 2}, \eqref{eq: sum Nj condition} and the fact that $S_{i+1} \geq S_i,$ we have that $$\tfrac{3S_{i}}{2} \leq \tfrac{3S_{i+1}}{2} \leq R_i \leq 6S_{i+1}, \qquad i \in \N.$$ Together with \eqref{eq: Si 2}, we have that
  \begin{equation} \label{eq: niriomegai}
\frac{N_{i+1}}{R_i\omega\left(\frac{1}{R_i}\right)} \geq \frac{N_{i+1}}{6S_{i+1}\omega\left(\frac{2}{3S_{i}}\right)}\geq \frac{N_{i+1}}{6S_{i+1}\omega\left(\frac{1}{S_{i}}\right)}\geq \frac{1}{6}, \qquad i \in \N. \end{equation}
 Noting by the definition of $g^{(d)}$ that $(g^{(d)})^{-1}(x_i) = x_i+(2N_{i+1},0,...,0)$ and $(g^{(d)})^{-1}(y_i) = y_i + (2,0,...,0) = x_i-(1,0,...,0),$ we have that
  \begin{equation*}
      \begin{split}
          \frac{||(g^{(d)})^{-1}(x_i) - (g^{(d)})^{-1}(y_i)||}{R_i \widetilde{\omega}\left(\frac{1}{R_i}\right)}  &= \frac{2N_{i+1}+1}{R_i \widetilde{\omega}\left(\frac{1}{R_i}\right)} \geq \frac{N_{i+1}}{R_i \widetilde{\omega}\left(\frac{1}{R_i}\right)}
          \geq \frac{1}{6}\frac{{\omega}\left(\frac{1}{R_i}\right)}{\widetilde{\omega}\left(\frac{1}{R_i}\right)}\xrightarrow{i \to \infty}\infty,
      \end{split}
  \end{equation*} where for the final inequality we used \eqref{eq: niriomegai}. Hence, $g^{(d)}$ is not $\widetilde{\omega}-$co-homogeneous.
\end{proof}

\end{document}
